\section{Reading the Outputs}
\label{sec:outputs}

\subsection{The output folder}

\begin{lstlisting}[language={}]
output/<study>/
|-- summary.md          the whole sweep in one table, robustness, heating, design
|-- run.log             what the run did (same as the console)
|-- config_used.yaml    every value used, defaults included
|-- tables/             CSV/JSON/TXT results (below)
|-- figures/            one PDF + PNG per plot
|-- cad/                profile CSV, STL and STEP files
|-- latex/              LaTeX tables and macros (if output.latex)
\end{lstlisting}

\subsection{Tables}

\paragraph{\texttt{trade\_table.csv}: the sweep.} One row per nose (profile $\times$ fineness $\times$ tip radius),
the same content as the table of \texttt{summary.md} with full precision:

\begin{table}[H]
\centering
\footnotesize
\begin{tabular}{@{}>{\raggedright\arraybackslash}p{3.0cm}>{\raggedright\arraybackslash}p{13.0cm}@{}}
\toprule
Column & Meaning \\
\midrule
\texttt{shape}, \texttt{fineness} & profile name and $L/D$ (\texttt{tip\_radius\_ratio} is added if tip radii are swept) \\
\texttt{L\_mm}, \texttt{V\_L}, \texttt{S\_m2} & nose length [mm], internal volume [L], wetted area [m\textsuperscript{2}] \\
\texttt{mass\_g} & shell (wetted area $\times$ areal density) + tip insert [g] \\
\texttt{base\_angle} & angle between the profile and the body at the shoulder [\textdegree]; 0 = tangent, otherwise a kink \\
\texttt{xcp\_over\_L} & Barrowman centre of pressure of the nose, from the apex, as a fraction of $L$ \\
\texttt{l3U\_mm}, \texttt{bay\_mm} & length of the curved nose wide enough for the payload; nose-bay length (curved part + the cylinder the payload still needs, at least the shoulder) [mm] \\
\texttt{cd\_w\_11} \dots & nose wave-drag coefficient at $M = 1.1$, 1.2, 1.4, 1.6 (reference area $\pi D^2/4$) \\
\texttt{apogee}, \texttt{mach\_max} & apogee [m] and maximum Mach number of the nominal flight with this nose \\
\texttt{dSM\_cal} & change of static margin relative to the baseline nose [calibres]; negative = less stable \\
\texttt{d\_apogee}, \texttt{d\_bay\_mm} & apogee and nose-bay length relative to the baseline nose [m], [mm] \\
\bottomrule
\end{tabular}
\end{table}

\paragraph{Other tables.}
\begin{itemize}[leftmargin=*,itemsep=1pt]
\item \texttt{robustness.csv}: apogee of every profile in every scenario and \texttt{dz\_vs\_vk}, its difference
from the reference profile in the same scenario [m].
\item \texttt{catalogue.csv}: every profile at the design fineness (volume, area, angles, CP, wave drag computed
directly at that fineness when available, apogee changes, mass, payload room, static margin).
\item \texttt{trade\_matrix.csv}: scores (1--5) of each criterion and the weighted total.
\item \texttt{bluntness.csv}: wave-drag increase [\%] and apogee change [m] of each reported tip radius.
\item \texttt{heating.json}: peak temperatures [\textdegree C] (nominal, envelope) of the stagnation point, metal
tip, printed tip and shell inner face, and the maximum Mach numbers.
\item \texttt{design.csv}: geometry of the selected design: profile length $L$, physical length after blunting,
truncation, tangency point of the tip sphere, volume, area, CP, shoulder angle.
\item \texttt{calibration.txt}: the factor $k$; \texttt{sensitivity.txt}: apogee change per kg and per unit $C_D$.
\end{itemize}

\paragraph{CAD files (\texttt{cad/}).} Units are millimetres, $x$ from the (blunted) tip.
\texttt{nose\_profile\_final.csv}: the profile ($x$, $r$), points clustered at the tip, for CAD sketches or
templates. \texttt{nose\_<profile>\_LD<L/D>.stl}: every profile of the sweep at the design fineness;
\texttt{\dots\_FINAL.stl} and \texttt{\dots\_FINAL.step}: the selected design. All solids include the shoulder
and a flat base; the STL files are checked to be closed (no open edges) and the STEP solid is checked by reading
it back. STEP is the file to import in SpaceClaim/Fluent or CAD.

\paragraph{LaTeX (\texttt{latex/}).} \texttt{tab\_trade.tex}, \texttt{tab\_matrix.tex}, \texttt{tab\_catalogue.tex},
\texttt{tab\_check45.tex} and \texttt{macros.tex} (every number quoted in the design report, e.g.\
\verb|\FINAPO|). \verb|\input{}| them in a report so that the text always matches the results.

\subsection{Plots}
\label{sec:plots}

The figures below are those of the Weisshorn study. Every plot uses the same colour, line style and marker per
profile, so that profiles can be followed from one figure to the next, also in black and white.

\newcommand{\plot}[3]{\paragraph{\texttt{\detokenize{#1}}.} #3
\begin{figure}[H]\centering\includegraphics[width=#2\textwidth]{fig_#1}\end{figure}}

\plot{profiles}{0.95}{(a) Profiles at the baseline fineness, drawn to scale; the dotted line is the radius the
payload needs, so the part of each curve above it can house the payload. (b) Zoom on the tip with the selected
spherical blunting shaded. Use it to see how quickly each profile gains radius.}

\plot{schematics}{0.85}{How each profile is constructed (drawn at $L/D = 2$ so that the differences show):
upper label, the shoulder (tangent or kink); lower label, the generator and the tip; inset, the tip magnified.
Panel (a) defines the notation.}

\plot{cdw}{0.95}{(a) Nose wave drag versus Mach number at the baseline fineness: markers are CFD results, lines
the model used in the trajectory; in the grey band the model is an assumption. (b) Wave drag versus fineness at
$M = 1.4$ on logarithmic axes; parallel straight lines mean $C_{D,\rm w} \propto (L/D)^{-2}$, as slender-body theory
predicts (dotted line). Lower is better.}

\plot{cp}{0.95}{Wall pressure coefficient along the nose at two Mach numbers ($x/L = 1$: shoulder). Positive
$C_p$ pushes the nose back (drag), negative $C_p$ on the aft part pulls it forward. Explains \emph{why} a profile
has low drag.}

\plot{validation}{0.95}{Checks of the CFD solver. (a) Cone surface pressure from the solver (circles) against the
exact Taylor--Maccoll solution (lines): the circles must sit on the lines. (b) Wave drag on three grids divided by
the extrapolated value: the legend gives the observed order $p$ (should be close to 2) and the grid convergence
index (GCI, the estimated grid error of the production grid).}

\plot{trajectory}{0.95}{(a) Mach number versus time on the nominal and hot-day flights, with the maximum Mach
number; grey: transonic range. (b) Share of the momentum lost to drag in each Mach interval (light) and the part
due to the nose wave drag (dark). The nose shape only matters in the Mach range where the dark bars are.}

\plot{breakdown}{0.95}{Drag coefficient of the whole vehicle split into its components (coasting, at the altitude
given in the plot options). The nose wave drag is the top band; its share at $M = 1.4$ is annotated. This is why
nose shape changes the apogee by tens of metres, not kilometres.}

\plot{trade}{0.95}{The sweep on one figure: (a) apogee change relative to the baseline (cross), (b) length of the
nose able to house the payload (dotted: full payload length), (c) static-margin change, all versus fineness, one
line per profile. Look for the profile that is highest in (a) for an acceptable (b).}

\plot{robustness}{0.85}{Apogee of each profile relative to the reference profile under each modelling scenario
(one marker per scenario). If a profile stays on the same side of zero in every scenario, the ranking does not
depend on the uncertain assumptions. Grey band: differences too small to matter (\key{robustness.band}).}

\plot{bluntness}{0.95}{(a) Wave drag of the blunted nose relative to the sharp one versus tip radius, at three
Mach numbers (from the CFD bluntness study). (b) Stagnation-point heat-transfer coefficient at maximum Mach versus
tip radius ($\propto r_n^{-1/2}$). A larger tip lowers the heating but raises the drag; the dotted line is the
selected radius.}

\plot{thermal}{0.95}{(a) Inner-face temperature of the shell at the configured stations (solid: nominal flight,
dashed: hot day; black: stagnation temperature, the upper bound) with the adhesive limit. (b) Tip temperature for
a printed polymer tip and for the metal insert, with the polymer's heat-deflection range. Curves must stay below the
limits.}

\plot{final}{0.95}{The selected design to scale: profile, tip insert, shoulder and payload envelope, with the main
dimensions.}

\plot{gallery}{0.85}{Every profile at the design fineness, to scale, with the payload placed as far forward as
the profile allows (dashed box; dotted: extra cylinder needed) and its key numbers.}

\plot{fields_M14}{0.85}{Pressure coefficient around each profile from the saved CFD flow fields (red: compression
behind the shock, blue: expansion). Dashed: sonic line where a subsonic region exists.}

\plot{grid}{0.8}{The CFD grid (coarse level for clarity) and its boundaries; for explanations, not for decisions.}

\plot{renders}{0.8}{Shaded 3D views of every profile and of the selected design (tip insert in grey), for
presentations.}
